\section{Hipótesis}

Los grandes modelos de lenguaje (LLMs) no satisfacen dichos criterios. 
Lo que exhiben son \textit{inferencias inductivas} 
\parencite{fawkes2024fragility} y se encuentran en la primera capa de 
la jerarquía de causalidad de \textcite{pearl2018theoretical}, las 
\textit{asociaciones}. Un LLM que produce respuestas causalmente 
correctas sin tener acceso a los \textit{entrelazamientos} causales del sistema natural 
está, en términos de \textcite{rosen2012anticipatory}, 
imputando\footnote{Imputar, en el sentido de Rosen, significa atribuir 
al mundo externo una estructura relacional construida por el sistema 
observador (nuestra mente), tratándola como si fuera una propiedad objetiva de ese 
mundo y no una construcción del observador.} relaciones al mundo 
externo sin el respaldo de una relación de modelado legítima. No hay 
un sistema natural $N$ codificado en él de manera que sus inferencias 
se decodifiquen como predicciones verificables sobre ese $N$ 
específico. 

Esta distinción adquiere precisión formal a través de la noción de 
modelo causal de \textcite{pearl2000causality}. Para Pearl, un modelo 
causal no es simplemente una distribución de probabilidad sobre 
\textit{observables}, más bien es una representación explícita de los \textit{mecanismos} 
que se alteran cuando el mundo cambia. En sus propias palabras, los 
modelos causales \textcite[p.~203]{pearl2000causality} ``encode and distinguish information about external changes through an explicit representation of the mechanisms that are 
altered in such changes''. Esto es precisamente lo que Rosen exige de 
una relación de modelado legítima: que las inferencias del sistema 
formal correspondan a relaciones estructurales del sistema natural, no 
solo a regularidades observadas bajo condiciones estáticas.

Los LLMs no tienen acceso a esa representación mecanística. Operan 
sobre una distribución observacional comprimida en texto, sin la estructura que permitiría 
responder a preguntas sobre \textit{intervenciones} o \textit{contrafactuales}. Más aún, como señalan 
\textcite{scholkopf2021toward}, un modelo con comprensión estructural 
genuina requiere que sus mecanismos internos sean independientes entre 
sí, el principio de mecanismos causales independientes (ICM). Un LLM 
opera sobre una mezcla entrelazada (\textit{entangled}) de todos los 
patrones extraídos del corpus, no sobre la factorización causal 
(\textit{disentangled}) que el ICM requiere. La diferencia no es de grado sino de 
arquitectura; un científico que imputa una relación causal dispone de 
un modelo formal cuyos mecanismos son separables e intervenibles; un 
LLM dispone de una distribución de probabilidad sobre tokens optimizada 
para predecir texto humano.

\section{Resumen Hipotesis}